\section{Research objective and evidence boundary}
\label{sec:scope}

The objective is to make exact unknot recognition faster while moving toward
a defensible quasi-polynomial worst-case analysis. That requires a joint
argument about topology, representation size, primitive bit costs, and the
number of primitive calls. Speeding up one operation does not supply the
missing bounds on the other three. The work below therefore treats an
implemented theorem, an observed acceleration, and a proposed global route
as different kinds of evidence.

The baseline is the public ProveIt revision
\begin{center}\small\baseline\end{center}
dated 8 October 2026, 15:40:39 UTC. All code changes in the accompanying patch
are relative to that revision. The target is
\code{Topology/UnknotRecognition/fast}; the older reports and synthesis
documents provide mathematical context and selected independent reference
implementations. No remote repository changes are required to inspect or
reproduce this package.

\subsection{Input size and the asymptotic target}

An input knot diagram has $n$ crossings and a validated planar-diagram
encoding. Arc labels are normalized during validation; their accidental
decimal spelling is not a topological size parameter. We use bit complexity
throughout. Writing a bound as polynomial in $n$ permits the usual
$O(n\log(n+2))$ encoding overhead. For a triangulation, $T$ denotes the number
of tetrahedra and $B$ the bit length of supplied integral cocycle heights.
For a compressed word, the grammar size and the bit length of its represented
length are input parameters; the expanded length can be exponential in them.

For an input encoding of length $N$, quasi-polynomial time means
\[
 \exp\!\left((\log(N+2))^{O(1)}\right).
\]
In particular $N^{C\log N}$ is quasi-polynomial, whereas
$2^{C\sqrt N}$ is a larger asymptotic class. Removing a factor $\log n$ from
an exponent $O(\sqrt n\log n)$ is a real improvement, but it does not reach
the requested target. Conversely, a polynomial operation on a $2^N$-letter
word is efficient only if the polynomial is in the compressed encoding,
not in $2^N$.

Three separate resource notions recur:
\begin{enumerate}
\item \emph{Mathematical input size}: the diagram, triangulation, grammar,
and binary integers specified by the theorem.
\item \emph{Represented intermediate size}: frontier keys, coefficient
bits, grammar nodes, attaching maps, or normal coordinates actually stored.
\item \emph{Charged implementation work}: a documented counter or elapsed
time allowance. A cap on state count is not a byte-memory cap, and a cap on
work does not turn an unfinished search into a negative mathematical answer.
\end{enumerate}

\subsection{What the existing implementation already supplies}

The maintained recognizer is a portfolio. It validates the input, applies
structural and simplification certificates, uses one-sided invariant tests,
and can attempt independently replayed group certificates. Complete exact
fallbacks include the built-in Khovanov calculation and an optional external
normal-surface engine. The previous work already includes separator
certificates, exact faithful Potts evaluation, transfer methods for special
diagram families, compressed group words, accelerated repeated automorphisms,
whole-donor macros, and a complete local compressed substring primitive.

The relevant baseline guarantees are narrower than a complete recognition
bound. A crossing scan with certified edge frontier $O(\sqrt n)$ is already
available. The faithful Potts calculation uses equality partitions, whose
Bell-number state count gives $2^{O(w\log(w+2))}$ for $w$ active boundary
objects. Its universal full-polynomial bound is thus
$\poly(n)2^{O(\sqrt n\log n)}$. The compressed word kernel has polynomial
local bit bounds, but no general bound on discovery steps or accumulated
grammar size follows. The normal-surface adapter can return an exact decision
from its external engine, but it does not yet export an independently replayed
compressed hierarchy with a quasi-polynomial running-time theorem.

This continuation reuses those working parts. It changes the state space of
the full-Jones calculation, optimizes a concrete geometric primitive, and
reduces work spent finding a sufficient compressed overlap. The three routes
share a principle: a compact representation must preserve the information
needed by the next exact operation, and its certificate must describe that
operation's actual conclusion.

\subsection{Relationship to the literature}

The Kauffman state model and its normalization are classical
\cite{kauffman1987}. Ellenberg, Newman, Sawin, and Shi already proved a
polynomial-times-$2^{C\sqrt n}$ time and memory bound for the Kauffman bracket
using controlled tangle boundaries \cite{ellenberg2013}. General planar
Boolean tensor contraction and Jones tensor implementations provide further
context \cite{liuying2023,meichanetzidis2019,laakkonen2026}. Our contribution
here is a self-contained binary model that works with the repository's
existing arbitrary crossing cut order, its exact arithmetic, its input
validation, and its resource semantics. The square-root-exponential exponent
itself is established prior work.

The geometry route is motivated by Lackenby's hierarchy construction
\cite{lackenby2026}. It is necessary to distinguish the 2021 announcement of
a quasi-polynomial unknot-recognition algorithm \cite{oxford2021} from the
particular written argument audited here. Section 9 of the July 2026 preprint
does not specify all primitive running times; its iteration estimate depends
on maximal hierarchy length and maximal pattern-complexity. Section 10
improves the depth control, but the code in this package does not implement
all of that geometric machinery. Thus neither the announcement nor the
preprint is used as a black-box complexity proof for ProveIt's recognizer.

Integral optimization through network matrices is also classical. The
topological use of total unimodularity has substantial precedent, including
optimal homologous cycle algorithms \cite{dey2011}. We claim a precise
local-span optimization, certificate, and geometric preservation argument
for the implemented cocycle-seed family, without asserting literature-wide
priority for network integrality or for connected nonseparating surfaces.

\subsection{A current claim about Jones detection and its audited prerequisites}
\label{sec:jones-status}

The literature check was made on 8 October 2026. Carmi and Cohen's
June 2026 preprint explicitly claims that the Jones polynomial detects the
unknot \cite[Theorem~6.9]{carmicohen2026}. The arXiv record inspected on that
date lists version~1 only and displays no withdrawal notice. This claim
must be acknowledged; it cannot be resolved by repeating older descriptions
of the problem's status. Our theorems and recognition decisions do not assume
the claimed detection theorem.

There is a concrete reason for this independence. The pinned repository
already contains an audit of a prerequisite consequence of that preprint.
Lemma~5.8, combined with the universal realization asserted in
Proposition~5.12 and the threshold in Section~5.1.1, implies a uniform bound
on the degree in $p=x-x^{-1}$ of every ordinary residual twist state:
\[
 \deg_p V\le B_{\rm host}+11+2.
\]
Here the expansion is $V(x)=a(p)+b(p)x$, and the degree means
$\max(\deg a,\deg b)$. This implication precedes the low-order anchor
vanishing used later in Theorem~5.10; it is not restricted to hosts already
known to have trivial Jones polynomial \cite[Sections~5.1--5.6]{carmicohen2026}.

For clarity, the existing witness can be described without any clasp-machine
isotopy. Take the three-crossing host with planar-diagram rows
\[
 (5,1,2,0),\quad(1,4,3,2),\quad(0,3,4,5).
\]
Replace the third row by an odd anti-parallel twist of length $L$, keeping
its exterior ports fixed. Explicitly, the first twist row begins with
$(a_0,b_0)=(0,3)$. At row $i=0,\ldots,L-1$ use
$(a_i,b_i,c_i,d_i)$, where $(c_i,d_i)=(6+2i,7+2i)$ for $i<L-1$,
$(c_{L-1},d_{L-1})=(4,5)$, and
$(a_{i+1},b_{i+1})=(d_i,c_i)$. These are the same fixed-exterior
diagrams used by the baseline audit.

Their knot Jones polynomials, in $t=x^2$, satisfy the audited skein recurrence
\[
 V_1(t)=t+t^3-t^4,\qquad
 V_L(t)=t^2V_{L-2}(t)+t-t^2+t^3-t^4.
\]
Consequently the highest term is $-t^{L+3}$. The identity
$x^m=F_{m-1}(p)+F_m(p)x$, with
\[
 F_m(p)=\sum_{k=0}^{\lfloor(m-1)/2\rfloor}
          \binom{m-k-1}{k}p^{m-2k-1},
\]
shows that the corresponding $b$-polynomial has leading term
$-p^{2L+5}$. All lower $t$-powers contribute smaller $p$-degrees.
Thus the degree is $2L+5$; for the mirror it is $2L+6$.

The host has $B_{\rm host}=3\cdot3+1-1=9$, so the proposed uniform
residual cap is $22$. The baseline's $L=9$ witness has degree $23$.
The present package independently recomputes its complete Jones polynomial
with Regina~7.4.1, and recomputes the formal expansion using the displayed
binomial formula rather than the baseline's coefficient recurrence.
It also checks $L=23$: the degree is $51$, while $L>22$ even satisfies
the long-twist inequality needed by the later descent. The mirror degrees
are $24$ and $52$, respectively. All four polynomial computations agree
coefficient by coefficient with all three tensor arithmetic modes.

The original audit additionally records local oriented-smoothing reductions
and Reidemeister-II anchor cancellations. Our new independent computation
replicates its polynomial obstruction. It neither constructs a nontrivial
knot with Jones polynomial one nor identifies a unique erroneous
clasp-machine step. It shows that the stated universal geometric realization
and uniform residual-degree consequence cannot both apply to these explicit
ordinary twists. Accordingly, a Jones polynomial different from one gives
the usual nontriviality obstruction, whereas identity remains inconclusive
in this implementation. The source and complete raw outputs are
\path{reproduce/audit_regina_jones.py} and
\path{results/regina_jones_20261008.json}.


\subsection{The principal results and their scope}

\begin{table}[H]
\centering\small
\begin{tabularx}{\linewidth}{@{}P{0.20\linewidth}YY@{}}
\toprule
Primitive & Proved improvement & Remaining limitation\\
\midrule
Full Jones polynomial & Binary frontier states; exact
$\poly(n)2^{O(\sqrt n)}$ bit bound with a certified separator order &
Subexponential rather than quasi-polynomial; identity is not accepted as an
unknot certificate by this implementation.\\[4pt]
Cocycle normal seed & Exact minimum disc count within a fixed
vertex-potential family; polynomial binary cost and a short dual certificate &
Restricted family; no incompressibility, genus-minimality, or hierarchy-reset
bound.\\[4pt]
Compressed overlap & A sufficient strict-majority witness followed by at most
two LCP extensions; exact local-existence completeness when uncapped &
May choose a worse greedy move; preceding stages and total search remain
possible bottlenecks.\\
\bottomrule
\end{tabularx}
\caption{Theorem scope. Each result is used only for the conclusion it proves.}
\label{tab:scope}
\end{table}

The paper's algorithmic claims are proved below and checked against independent
oracles where feasible. Those checks increase confidence in the implementation;
finite tests do not establish the asymptotic theorems. Conversely, a correct
proof of an algorithm does not excuse a mismatch between the mathematical
algorithm and the code. The package therefore retains full inputs, exact
outputs, proof traces where available, source hashes, and censored observations.
